% GENERATED_BY: oc_core_monograph_translation_draft_pipeline_v1.py
% AI_ASSISTED_TRANSLATION_DRAFT: TRUE
% THEOREM_REVIEW_STATUS: PENDING
% THEOREM_REVIEW_MODE: FORMAL_AI_GATE__CODEX_CLI_CHATGPT
% THEOREM_REVIEW_REVIEWER_ID:
% THEOREM_REVIEW_AUTHORITY_CLASS:
% THEOREM_REVIEWED_AT:
% SOURCE_LANGUAGE: EN
% TARGET_LANGUAGE: RU
% SOURCE_EN_REF: content/m_spaces/m5.tex
% ================================================================
% ==== FILE: content/m_spaces/m5.tex
% ================================================================

% ==============================
%  Ontology of Continua — Core
%  Meta-Space: M5
%  FULL MODULE — FINAL VERSION
% ==============================

\subsubsection{\texorpdfstring{$M_5$}{M_5} Обзор}
\label{sec:m5-overview}

$M_5$ — это мета-пространство, которое делает \emph{электрическую возбудимость} допустимой.
Это среда, в которой биологический континуум может проявлять:

\begin{itemize}
    \item стабильный мембранный потенциал $V(t)$,
    \item проводимостьные состояния, опосредованные ионными каналами,
    \item возбуждение на основе порога ($\Theta_{\mathrm{exc}}$),
    \item распространяющиеся электрические события (предшественники потенциалов действия),
    \item регенеративную обратную связь, необходимую для спайк-циклов.
\end{itemize}

Где $M_4$ обеспечивало \emph{химическую} компартментализацию,
$M_5$ обеспечивает \emph{электрохимическую} динамику, которая определяет континуум $K_5$.

% ---------------------------------------------------------------
\subsubsection*{1. Допустимое пространство конфигураций \texorpdfstring{$\Omega(M_5)$}{\Omega(M_5)}}

Допустимые конфигурации расширяют $\Omega(M_4)$, включая:

\[
\Omega(M_5) =
\left\{
(\partial\Omega,\ V,\ g_{\mathrm{ion}},\
\Pi_{\mathrm{ion}},\
\Delta\mu_{\mathrm{ion}},\
C_m,\
\mathcal{C}_{\mathrm{channels}})
\ \Big|\
\text{выполняются условия допустимости}
\right\}.
\]

Здесь:

\begin{itemize}
    \item $V$ — поле мембранного потенциала,
    \item $C_m$ — ёмкость мембраны,
    \item $g_{\mathrm{ion}}$ — состояния проводимости ионных каналов,
    \item $\mathcal{C}_{\mathrm{channels}}$ — популяция каналов и логика их гейтинга,
    \item $\Pi_{\mathrm{ion}}$ — ион-специфические проницаемости,
    \item $\Delta\mu_{\mathrm{ion}}$ — электрохимические градиенты, порождающие токи.
\end{itemize}

Допустимость требует:

\begin{enumerate}
    \item Целостность мембраны, как в $M_4$.
    \item Ионные каналы, способные переключаться минимум между двумя состояниями:
          \[
          \{ \text{открыт},\ \text{закрыт} \}.
          \]
    \item Динамика потенциала, удовлетворяющая:
          \[
          C_m \frac{dV}{dt} = - \sum_i g_i (V - E_i) + I_{\mathrm{ext}}
          \]
          в области допустимых параметров.
    \item Наличие нетривиального равновесия и сепаратрисы, определяющей порог.
\end{enumerate}

Следовательно, $M_5$ — минимальная среда, где клетка может быть возбудимой.

% ---------------------------------------------------------------
\subsubsection*{2. Допустимые оси \texorpdfstring{$A(M_5)$}{A(M_5)}}

Возникает новая фундаментальная ось:

\[
A_{\mathrm{exc}} = \{ \text{возбудимая},\ \text{невозбудимая} \},
\]

с соответствующим различием по пороговому напряжению:

\[
A_V = \{ V < \Theta_{\mathrm{exc}},\ V \ge \Theta_{\mathrm{exc}} \}.
\]

Другие допустимые оси:

\begin{itemize}
    \item $A_{\mathrm{ion}}$ — ион-специфические конформации каналов,
    \item $A_{\mathrm{gate}}$ — состояния гейтинговых переменных (активация/инактивация),
    \item $A_{\mathrm{cond}}$ — режимы проводимости,
    \item $A_{\mathrm{spike}}$ — фазы рефрактерности/перерегулирования/сброса.
\end{itemize}

Эти оси определяют высокоразмерную динамическую геометрию возбудимых мембран.

% ---------------------------------------------------------------
\subsubsection*{3. Потенциалы \texorpdfstring{$P(M_5)$}{P(M_5)}}

$M_5$ вводит потенциалы, связанные с электро-динамикой:

\[
P(M_5) =
(U_{\mathrm{ion}},\ U_V,\ U_{\mathrm{gate}},\
U_{\mathrm{CAP}},\ U_{\mathrm{rest}}).
\]

Интерпретация:

\begin{itemize}
    \item $U_{\mathrm{ion}}$ — электрохимические потенциалы, подобные потенциалу Нернста,
    \item $U_V$ — энергетический ландшафт, зависящий от напряжения,
    \item $U_{\mathrm{gate}}$ — энергия гейтинга для переходов каналов,
    \item $U_{\mathrm{CAP}}$ — ёмкостная энергия мембраны,
    \item $U_{\mathrm{rest}}$ — минимум энергии состояния покоя.
\end{itemize}

Условия допустимости:

\begin{enumerate}
    \item $U_V$ должен иметь локальный минимум (состояние покоя),
    \item функция усиления должна допускать сверхпороговую нестабильность,
    \item переходы гейтинга должны быть термически и химически осуществимыми.
\end{enumerate}

% ---------------------------------------------------------------
\subsubsection*{4. Пороги \texorpdfstring{$\Theta(M_5)$}{\Theta(M_5)}}

Определяющий порог:

\[
\Theta_{\mathrm{exc}} = \text{минимальная деполяризация, необходимая для запуска регенеративного события}.
\]

Дополнительные пороги:

\begin{itemize}
    \item $\Theta_{\mathrm{gate}}$ — порог активации гейтинга,
    \item $\Theta_{\mathrm{cond}}$ — порог переключения проводимости,
    \item $\Theta_{\mathrm{stability}}$ — граница устойчивости для осцилляторных режимов.
\end{itemize}

Критическое соотношение:

\[
V(t) \ge \Theta_{\mathrm{exc}} \quad\Rightarrow\quad
\text{переход в режим генерации спайка}.
\]

% ---------------------------------------------------------------
\subsubsection*{5. Потоки \texorpdfstring{$J(M_5)$}{J(M_5)}}

Допустимые потоки расширяют потоки из $M_4$, включая:

\[
J(M_5) =
\big(
J_{\mathrm{ion}},\
J_V,\
J_{\mathrm{gate}},\
J_{\mathrm{exc}},\
\nabla_i T_5
\big).
\]

Здесь:

\begin{itemize}
    \item $J_{\mathrm{ion}}$ — ион-специфические токи,
    \item $J_V$ — потенциал-зависимые потоки,
    \item $J_{\mathrm{gate}}$ — динамические переходы гейтинга,
    \item $J_{\mathrm{exc}}$ — режим регенеративной возбуждающей динамики,
    \item $\nabla_i T_5$ — потоки, вызванные градиентами напряжения при возбудимости.
\end{itemize}

Законы сохранения:

\[
\sum_i J_{\mathrm{ion},i} = C_m \frac{dV}{dt}.
\]

% ---------------------------------------------------------------
\subsubsection*{6. Циклы в \texorpdfstring{$M_5$}{M_5}}

Ключевой цикл:

\[
C_{\mathrm{spike}}
=
\{\text{покой} \to \text{деполяризация} \to
\text{переразгон} \to \text{реполяризация} \to \text{рефрактерность} \to \text{покой}\}.
\]

Другие допустимые циклы:

\begin{itemize}
    \item $C_{\mathrm{gate}}$ — петля активации/инактивации гейта,
    \item $C_{\mathrm{osc}}$ — субпороговые осцилляторные циклы (если разрешены),
    \item $C_{\mathrm{conduct}}$ — переходные циклы между состояниями проводимости.
\end{itemize}

Условие существования спайк-цикла:

\[
\oint_{C_{\mathrm{spike}}} dA_{\mathrm{phase}} \approx 0
\quad\text{и}\quad
V(t) \text{ пересекает } \Theta_{\mathrm{exc}} \text{ один раз за цикл}.
\]

% ---------------------------------------------------------------
\subsubsection*{7. Граница \texorpdfstring{$\partial\Omega(M_5)$}{\partial\Omega(M_5)}}

Конфигурация приближается к $\partial\Omega(M_5)$, когда:

\begin{itemize}
    \item мембрана не может поддерживать стабильный потенциал покоя,
    \item кинетика гейтинга нарушает возбудимость,
    \item проводимость насыщается или коллапсирует,
    \item $\Theta_{\mathrm{exc}}$ становится недостижимой (слишком высокой) или тривиальной (нуль),
    \item персистирует деполяризационный блок,
    \item ионные градиенты падают ниже минимально жизнеспособных уровней.
\end{itemize}

Пересечение $\partial\Omega(M_5)$ делает спайк невозможным;
поэтому $K_5$ коллапсирует или возвращается к режиму, похожему на $K_4$.

% ---------------------------------------------------------------
\subsubsection*{8. Связь с \texorpdfstring{$M_4$}{M_4} и \texorpdfstring{$M_6$}{M_6}}

Переход $K_4 \to K_5$ допустим, когда:

\[
A_{\mathrm{exc}} \in A(M_5),
\qquad
V(t) \ge \Theta_{\mathrm{exc}},
\qquad
\text{ионные каналы формируют когерентную динамику гейтинга}.
\]

Переход к $K_6$ (когнитивные аттракторы) требует:

\[
A_{\mathrm{concept}} \in A(M_6),
\qquad
C_{\mathrm{spike}} \text{ должен поддерживать сопряжение с другими спайками},
\]
\[
\text{возникновение аттракторной динамики из спайк-трендов}.
\]

Следовательно, $M_5$ — это мета-пространство, обеспечивающее:

\begin{itemize}
    \item электрическую возбудимость,
    \item генерацию спайков,
    \item динамическую регуляцию проводимости,
    \item раннюю биоэлектрическую логику,
    \item структурный мост к нейронным вычислениям.
\end{itemize}
